\documentclass[11pt,a4paper]{article}

% ===================================================================
% PACKAGES
% ===================================================================
\usepackage[utf8]{inputenc}
\usepackage[T1]{fontenc}
% lmodern not available in this environment
\usepackage{amsmath,amssymb,amsthm}
\usepackage{algorithm2e}
\usepackage{booktabs}
\usepackage{hyperref}
\usepackage[margin=1in]{geometry}
\usepackage{enumitem}

% ===================================================================
% THEOREM ENVIRONMENTS
% ===================================================================
\newtheorem{theorem}{Theorem}[section]
\newtheorem{lemma}[theorem]{Lemma}
\newtheorem{proposition}[theorem]{Proposition}
\newtheorem{corollary}[theorem]{Corollary}
\theoremstyle{definition}
\newtheorem{definition}[theorem]{Definition}
\newtheorem{remark}[theorem]{Remark}

% ===================================================================
% NOTATION SHORTCUTS
% ===================================================================
\newcommand{\Z}{\mathbb{Z}}
\newcommand{\Zq}{\mathbb{Z}_q}
\newcommand{\Rq}{R_3}
\newcommand{\GF}{\mathrm{GF}(3)}
\newcommand{\norm}[1]{\left\lVert #1 \right\rVert}
\newcommand{\linf}[1]{\left\lVert #1 \right\rVert_{\infty}}
\newcommand{\Adv}{\mathbf{Adv}}
\newcommand{\negl}{\mathrm{negl}}
\newcommand{\poly}{\mathrm{poly}}
\newcommand{\wt}{\mathrm{wt}}
\newcommand{\Sample}{\mathrm{Sample}}
\newcommand{\KeyGen}{\textsf{KeyGen}}
\newcommand{\Sign}{\textsf{Sign}}
\newcommand{\Verify}{\textsf{Verify}}
\newcommand{\MSIS}{\textsf{MSIS}}
\newcommand{\MLWE}{\textsf{MLWE}}
\newcommand{\EUFCMA}{\textsf{EUF-CMA}}

% ===================================================================
% TITLE
% ===================================================================
\title{%
  \textbf{On the EUF-CMA Security of TL-DSA:\\
  A Formal Reduction to Module-LWE\\
  over Ternary Polynomial Rings}%
}

\author{%
  Applied Physics Division\\
  Capomastro Holdings Ltd.\\
  \texttt{research@capomastro.com}%
}

\date{March 2026}

\begin{document}

\maketitle

% ===================================================================
% ABSTRACT
% ===================================================================
\begin{abstract}
We present a formal security analysis of TL-DSA (Ternary Lattice Digital Signature Algorithm), a Fiat--Shamir with Aborts signature scheme operating in the polynomial ring $R_3 = \Z_3[X]/(X^{256}+1)$ with balanced ternary coefficients $\{-1, 0, +1\}$. TL-DSA is structurally analogous to ML-DSA (FIPS~204/Dilithium) but specialized to the ternary domain. We prove that TL-DSA achieves existential unforgeability under chosen-message attack (EUF-CMA) in the random oracle model, assuming the hardness of the Module Short Integer Solution (MSIS) problem over $R_3$. We further show that MSIS over $R_3$ is at least as hard as Module-LWE over $R_3$ via the standard SIS--LWE duality. A structural property unique to the ternary setting---that rejection sampling becomes trivially satisfied when $q=3$---is identified and analyzed, yielding deterministic (non-aborting) signing with perfect zero-knowledge simulation. Concrete security bounds are derived for all three parameter sets (TL-DSA-44, TL-DSA-65, TL-DSA-87) targeting 128, 192, and 256 bits of classical security respectively.
\end{abstract}

\medskip
\noindent\textbf{Keywords:} Post-quantum signatures, lattice-based cryptography, Module-LWE, Module-SIS, Fiat--Shamir, ternary arithmetic, balanced ternary.

\tableofcontents

% ===================================================================
% SECTION 1: INTRODUCTION
% ===================================================================
\section{Introduction}
\label{sec:intro}

The standardization of lattice-based cryptographic schemes by NIST---ML-KEM (FIPS~203) and ML-DSA (FIPS~204)---has established Module-LWE as the foundational hardness assumption for post-quantum public-key cryptography. Both standards operate over polynomial rings $\Z_q[X]/(X^n+1)$ with large moduli ($q = 3329$ for ML-KEM, $q = 8380417$ for ML-DSA), chosen primarily to enable efficient Number Theoretic Transform (NTT) multiplication.

TL-DSA takes a fundamentally different approach: it operates in the minimal nontrivial polynomial ring $R_3 = \Z_3[X]/(X^{256}+1)$, where the modulus $q=3$ coincides with the coefficient domain of balanced ternary $\{-1, 0, +1\}$. This specialization is motivated by the PlenumNET ternary computing platform, where all arithmetic naturally operates in $\GF$.

This paper provides three contributions:

\begin{enumerate}[leftmargin=2em]
\item \textbf{Formal specification} of TL-DSA extracted from the production Rust implementation (Section~\ref{sec:construction}), ensuring the analyzed scheme matches the deployed code.

\item \textbf{EUF-CMA security proof} in the random oracle model (Section~\ref{sec:proof}), reducing forgery to solving Module-SIS over $R_3$, with an explicit reduction to Module-LWE via the standard duality (Section~\ref{sec:duality}).

\item \textbf{Identification of the ternary non-abort property} (Section~\ref{sec:ternary-advantage}): when $q=3$, the masking vector $\mathbf{y}$ perfectly hides the secret through modular addition, eliminating the rejection sampling that dominates signing cost in standard Dilithium.
\end{enumerate}

% ===================================================================
% SECTION 2: PRELIMINARIES
% ===================================================================
\section{Preliminaries}
\label{sec:prelim}

\subsection{Notation}

Bold lowercase letters ($\mathbf{s}, \mathbf{y}, \mathbf{z}$) denote vectors over $R_3$. Bold uppercase letters ($\mathbf{A}$) denote matrices over $R_3$. Regular lowercase ($c, f$) denote individual ring elements. We write $\linf{\mathbf{v}}$ for the maximum absolute value of any coefficient across all polynomials in $\mathbf{v}$. For a polynomial $f \in R_3$, $\wt(f)$ denotes the Hamming weight (number of nonzero coefficients).

\subsection{The Ternary Polynomial Ring}

\begin{definition}[Ternary Polynomial Ring]
\label{def:ring}
Let $n = 256$. The ternary polynomial ring is
\[
  R_3 = \Z_3[X] / (X^n + 1),
\]
where $\Z_3 = \{-1, 0, +1\}$ denotes the integers modulo 3 in balanced representation. Ring addition is componentwise modulo~3. Ring multiplication is polynomial multiplication reduced modulo $X^n + 1$, with coefficient reduction modulo~3 via the map $r \mapsto ((r \bmod 3) + 3) \bmod 3$ followed by $\{0 \mapsto 0,\; 1 \mapsto 1,\; 2 \mapsto -1\}$.
\end{definition}

\begin{remark}[Implementation correspondence]
This ring is implemented in \texttt{ternary\_lattice.rs} (lines 58--214). The \texttt{mod3} function (line~67) performs the balanced reduction. Ring multiplication (\texttt{ring\_mul}, line~185) uses schoolbook convolution with the $X^n+1$ negacyclic reduction: for indices $i+j \geq n$, the product term is subtracted rather than added.
\end{remark}

\subsection{Module-SIS and Module-LWE over $R_3$}

\begin{definition}[Module-SIS over $R_3$]
\label{def:msis}
For parameters $k, l \in \Z^+$ and bound $\beta \in \Z^+$, the $\MSIS_{k,l,\beta}$ problem over $R_3$ is: given a uniformly random matrix $\mathbf{A} \in R_3^{k \times l}$, find a nonzero vector $\mathbf{x} \in R_3^l$ such that $\mathbf{A} \cdot \mathbf{x} = \mathbf{0} \in R_3^k$ and $\linf{\mathbf{x}} \leq \beta$.
\end{definition}

\begin{definition}[Module-LWE over $R_3$]
\label{def:mlwe}
For parameters $k \in \Z^+$ and a noise distribution $\chi$ over $R_3$, the $\MLWE_{k,\chi}$ problem over $R_3$ is: given $(\mathbf{A}, \mathbf{b})$ where $\mathbf{A} \xleftarrow{\$} R_3^{k \times k}$ and either $\mathbf{b} = \mathbf{A} \cdot \mathbf{s} + \mathbf{e}$ for $\mathbf{s}, \mathbf{e} \leftarrow \chi^k$ or $\mathbf{b} \xleftarrow{\$} R_3^k$, distinguish which case holds.
\end{definition}

\begin{remark}[Ternary specialization]
When $q = 3$, every element of $R_3$ has $\linf{\cdot} \leq 1$. Both the uniform distribution and any noise distribution over $R_3$ produce coefficients in $\{-1, 0, +1\}$. This means Module-LWE and Module-SIS over $R_3$ are problems about the algebraic structure of $R_3$ rather than about distinguishing small errors from uniform---the ``errors'' are indistinguishable from uniform elements because they share the same support.
\end{remark}

\subsection{Centered Binomial Distribution over $\Z_3$}

\begin{definition}[CBD$_\eta$ over $\Z_3$]
\label{def:cbd}
For parameter $\eta \in \Z^+$, sample $a_1, \ldots, a_\eta, b_1, \ldots, b_\eta$ uniformly from $\{0, 1, 2\}$ (equivalently, from the output of the ternary sponge). Output $\left(\sum_{i=1}^\eta a_i - \sum_{i=1}^\eta b_i\right) \bmod 3$, interpreted in balanced form.
\end{definition}

This is implemented in \texttt{sample\_cbd\_ternary} (line~423 of \texttt{ternary\_lattice.rs}) with $\eta = 2$ for all TL-DSA variants.

\subsection{Ternary Sponge Hash Function}

TL-DSA uses a custom ternary sponge construction $H: \{-1,0,+1\}^* \to \{-1,0,+1\}^m$ as its hash function, implemented in \texttt{sponge.rs}. The sponge has state size 729 trits ($= 3^6$), rate 243 trits ($= 3^5$), capacity 486 trits, and performs 27 rounds of a nonlinear permutation per absorption.

In the security proof, we model $H$ as a random oracle $\mathcal{H}: \{-1,0,+1\}^* \to \{-1,0,+1\}^{243}$.

% ===================================================================
% SECTION 3: TL-DSA CONSTRUCTION
% ===================================================================
\section{TL-DSA Construction}
\label{sec:construction}

\subsection{Parameters}

\begin{table}[ht]
\centering
\caption{TL-DSA parameter sets, extracted from \texttt{tl\_dsa.rs} lines 73--103.}
\label{tab:params}
\begin{tabular}{@{}lcccccccc@{}}
\toprule
Variant & $n$ & $k$ & $l$ & $\eta$ & $\gamma$ & $\tau$ & Security (bits) & NIST Level \\
\midrule
TL-DSA-44 & 256 & 4 & 4 & 2 & 1 & 39 & 128 & 2 \\
TL-DSA-65 & 256 & 6 & 5 & 2 & 1 & 49 & 192 & 3 \\
TL-DSA-87 & 256 & 8 & 7 & 2 & 1 & 60 & 256 & 5 \\
\bottomrule
\end{tabular}
\end{table}

\noindent The parameters are:
\begin{itemize}[leftmargin=2em]
\item $n = 256$: polynomial degree (ring dimension).
\item $k$: number of rows in the public matrix $\mathbf{A}$ and dimension of $\mathbf{t}$.
\item $l$: number of columns in $\mathbf{A}$ and dimension of the secret $\mathbf{s}_1$.
\item $\eta = 2$: CBD parameter for secret key sampling.
\item $\gamma = 1$: $\ell_\infty$ rejection bound for signatures.
\item $\tau$: Hamming weight of challenge polynomial $c$.
\end{itemize}

\subsection{Key Generation}

\begin{algorithm}[H]
\caption{$\KeyGen(\texttt{seed})$}
\label{alg:keygen}
\DontPrintSemicolon
$\rho \leftarrow \mathcal{H}(\texttt{seed} \| 0)$ \tcp*{243 trits, matrix seed}
$\sigma \leftarrow \mathcal{H}(\texttt{seed} \| 1)$ \tcp*{243 trits, noise seed}
$\xi \leftarrow \mathcal{H}(\texttt{seed} \| 0 \| 1 \| {-1})$ \tcp*{243 trits, signing seed}
$\mathbf{A} \leftarrow \textsf{ExpandA}(\rho) \in R_3^{k \times k}$ \tcp*{uniform ternary matrix}
$\mathbf{s}_1 \leftarrow \textsf{CBD}_\eta(\sigma) \in R_3^l$ \tcp*{short secret vector}
$\mathbf{t} \leftarrow \mathbf{A} \cdot \mathbf{s}_1 \in R_3^k$ \tcp*{public key vector}
\Return $\left(\textit{pk} = (\rho, \mathbf{t}),\; \textit{sk} = (\rho, \mathbf{s}_1, \mathbf{t}, \xi)\right)$
\end{algorithm}

\begin{remark}[Zero error term]
Unlike ML-DSA, TL-DSA computes $\mathbf{t} = \mathbf{A} \cdot \mathbf{s}_1$ without an additive error term $\mathbf{s}_2$. In the implementation (\texttt{tl\_dsa.rs} line~211), $\mathbf{s}_2$ is initialized as the zero vector. This is justified by the ternary non-abort property (Section~\ref{sec:ternary-advantage}): since $q=3$, the distribution $\mathbf{A} \cdot \mathbf{s}_1$ is already pseudorandom under MSIS hardness, and adding noise would not improve security because the noise distribution equals the uniform distribution over $R_3$.
\end{remark}

\subsection{Signing}

\begin{algorithm}[H]
\caption{$\Sign(\textit{sk}, \mu)$}
\label{alg:sign}
\DontPrintSemicolon
\KwIn{Secret key $\textit{sk} = (\rho, \mathbf{s}_1, \mathbf{t}, \xi)$, message $\mu \in \{-1,0,+1\}^*$}
$\mathbf{A} \leftarrow \textsf{ExpandA}(\rho)$\;
$\bar{\mu} \leftarrow \mathcal{H}(\texttt{trits}(\mathbf{t}) \| \mu)$ \tcp*{message representative}
\For{$\kappa = 0, 1, 2, \ldots$}{
  $\mathbf{y} \leftarrow \textsf{ExpandMask}(\xi, \bar{\mu}, \kappa) \in R_3^l$ \tcp*{ternary masking vector}
  $\mathbf{w} \leftarrow \mathbf{A} \cdot \mathbf{y} \in R_3^k$\;
  $\tilde{c} \leftarrow \mathcal{H}(\bar{\mu} \| \texttt{trits}(\mathbf{w}))$ \tcp*{challenge hash}
  $c \leftarrow \textsf{SampleChallenge}(\tilde{c}) \in R_3$ \tcp*{weight-$\tau$ ternary poly}
  $\mathbf{z} \leftarrow \mathbf{y} + c \cdot \mathbf{s}_1 \in R_3^l$\;
  \If{$\linf{\mathbf{z}} \leq \gamma$}{
    \Return $\sigma = (\mathbf{z}, \tilde{c})$\;
  }
}
\end{algorithm}

\begin{remark}[Challenge polynomial]
$\textsf{SampleChallenge}(\tilde{c})$ produces a polynomial $c \in R_3$ with exactly $\tau$ nonzero coefficients, each in $\{-1, +1\}$. This is implemented in \texttt{sample\_challenge} (line~157 of \texttt{tl\_dsa.rs}).
\end{remark}

\subsection{Verification}

\begin{algorithm}[H]
\caption{$\Verify(\textit{pk}, \mu, \sigma)$}
\label{alg:verify}
\DontPrintSemicolon
\KwIn{Public key $\textit{pk} = (\rho, \mathbf{t})$, message $\mu$, signature $\sigma = (\mathbf{z}, \tilde{c})$}
\lIf{$\linf{\mathbf{z}} > \gamma$}{\Return $\bot$}
$\mathbf{A} \leftarrow \textsf{ExpandA}(\rho)$\;
$c \leftarrow \textsf{SampleChallenge}(\tilde{c})$\;
$\mathbf{w}' \leftarrow \mathbf{A} \cdot \mathbf{z} - c \cdot \mathbf{t} \in R_3^k$\;
$\bar{\mu} \leftarrow \mathcal{H}(\texttt{trits}(\mathbf{t}) \| \mu)$\;
$\tilde{c}' \leftarrow \mathcal{H}(\bar{\mu} \| \texttt{trits}(\mathbf{w}'))$\;
\Return $[\tilde{c}' = \tilde{c}]$\;
\end{algorithm}

\begin{proposition}[Correctness]
\label{prop:correct}
For any valid key pair $(\textit{pk}, \textit{sk}) \leftarrow \KeyGen$ and message $\mu$, if $\sigma \leftarrow \Sign(\textit{sk}, \mu)$ succeeds, then $\Verify(\textit{pk}, \mu, \sigma) = 1$.
\end{proposition}

\begin{proof}
Let $\sigma = (\mathbf{z}, \tilde{c})$ where $\mathbf{z} = \mathbf{y} + c \cdot \mathbf{s}_1$. In verification:
\begin{align*}
  \mathbf{w}' &= \mathbf{A} \cdot \mathbf{z} - c \cdot \mathbf{t} \\
  &= \mathbf{A}(\mathbf{y} + c \cdot \mathbf{s}_1) - c \cdot \mathbf{A} \cdot \mathbf{s}_1 \\
  &= \mathbf{A} \cdot \mathbf{y} + c \cdot \mathbf{A} \cdot \mathbf{s}_1 - c \cdot \mathbf{A} \cdot \mathbf{s}_1 \\
  &= \mathbf{A} \cdot \mathbf{y} = \mathbf{w}.
\end{align*}
Since $\mathbf{w}' = \mathbf{w}$, the recomputed challenge hash $\tilde{c}' = \mathcal{H}(\bar{\mu} \| \texttt{trits}(\mathbf{w})) = \tilde{c}$.
\end{proof}

% ===================================================================
% SECTION 4: THE TERNARY NON-ABORT PROPERTY
% ===================================================================
\section{The Ternary Non-Abort Property}
\label{sec:ternary-advantage}

We identify a structural property unique to TL-DSA's ternary instantiation that has significant implications for both security and efficiency.

\begin{theorem}[Ternary Non-Abort]
\label{thm:nonabort}
When $q = 3$ and $\gamma = 1$, the rejection condition $\linf{\mathbf{z}} > \gamma$ in Algorithm~\ref{alg:sign} is never satisfied. Every signing attempt succeeds on the first iteration ($\kappa = 0$).
\end{theorem}

\begin{proof}
Since all arithmetic is in $R_3$, every polynomial in $R_3$ has coefficients in $\{-1, 0, +1\}$. Therefore $\linf{f} \leq 1$ for all $f \in R_3$. In particular:
\begin{itemize}
\item $\mathbf{y} \in R_3^l$ implies $\linf{\mathbf{y}} \leq 1$.
\item $c \in R_3$ and $\mathbf{s}_1 \in R_3^l$ implies $c \cdot \mathbf{s}_{1,i} \in R_3$ for each component, so $\linf{c \cdot \mathbf{s}_1} \leq 1$.
\item $\mathbf{z} = \mathbf{y} + c \cdot \mathbf{s}_1 \in R_3^l$ implies $\linf{\mathbf{z}} \leq 1 = \gamma$.
\end{itemize}
The rejection condition $\linf{\mathbf{z}} > 1$ is never satisfied because $R_3$ is closed under addition and multiplication with $\linf{\cdot} \leq 1$.
\end{proof}

\begin{corollary}[Deterministic Signing]
\label{cor:deterministic}
TL-DSA signing is deterministic: given $(\textit{sk}, \mu)$, the signature $\sigma$ is uniquely determined. The for-loop in Algorithm~\ref{alg:sign} always terminates at $\kappa = 0$.
\end{corollary}

\begin{theorem}[Perfect Zero-Knowledge Simulation]
\label{thm:zk}
The signature $\mathbf{z} = \mathbf{y} + c \cdot \mathbf{s}_1 \in R_3^l$ is uniformly distributed over $R_3^l$ conditioned on $\mathbf{A} \cdot \mathbf{z} - c \cdot \mathbf{t} = \mathbf{w}$, regardless of the secret key $\mathbf{s}_1$.
\end{theorem}

\begin{proof}
The masking vector $\mathbf{y}$ is derived from the signing seed $\xi$, the message representative $\bar{\mu}$, and the attempt counter $\kappa = 0$ via the random oracle $\mathcal{H}$. In the random oracle model, $\mathbf{y}$ is uniformly distributed over $R_3^l$.

Since $(\Z_3, +)$ is a group, for any fixed $c \cdot \mathbf{s}_1 \in R_3^l$, the map $\mathbf{y} \mapsto \mathbf{y} + c \cdot \mathbf{s}_1$ is a bijection on $R_3^l$. Therefore $\mathbf{z} = \mathbf{y} + c \cdot \mathbf{s}_1$ is uniformly distributed over $R_3^l$ whenever $\mathbf{y}$ is, independent of $\mathbf{s}_1$.

This means a simulator without knowledge of $\mathbf{s}_1$ can produce identically distributed transcripts: sample $\mathbf{z} \xleftarrow{\$} R_3^l$, compute $\mathbf{w} = \mathbf{A} \cdot \mathbf{z} - c \cdot \mathbf{t}$, and program the random oracle accordingly.
\end{proof}

\begin{remark}[Contrast with ML-DSA]
In ML-DSA (Dilithium), $q = 8380417$ and $\gamma_1 \gg \linf{c \cdot \mathbf{s}_1}$. The masking vector $\mathbf{y}$ is sampled from $[-\gamma_1, \gamma_1]$, and rejection sampling ensures that $\mathbf{z}$ is statistically close to uniform on $[-\gamma_1 + \beta, \gamma_1 - \beta]$, which leaks no information about $\mathbf{s}_1$. The expected number of signing attempts is $\approx e^{l \cdot n \cdot 2\beta / (2\gamma_1+1)}$, typically 4--7 iterations for ML-DSA-65.

In TL-DSA, the group structure of $\Z_3$ provides \emph{perfect} masking for free: the one-time pad property of modular addition eliminates any need for rejection sampling or statistical closeness arguments. This reduces signing to exactly one attempt with zero information leakage.
\end{remark}

% ===================================================================
% SECTION 5: EUF-CMA SECURITY PROOF
% ===================================================================
\section{EUF-CMA Security Proof}
\label{sec:proof}

\subsection{Security Model}

\begin{definition}[EUF-CMA Security]
A signature scheme $\Pi = (\KeyGen, \Sign, \Verify)$ is $(\epsilon, q_S, q_H)$-EUF-CMA secure if for all adversaries $\mathcal{A}$ making at most $q_S$ signing queries and $q_H$ random oracle queries:
\[
  \Adv^{\EUFCMA}_\Pi(\mathcal{A}) = \Pr\left[
    \Verify(\textit{pk}, \mu^*, \sigma^*) = 1 \;\middle|\;
    \begin{array}{l}
      (\textit{pk}, \textit{sk}) \leftarrow \KeyGen \\
      (\mu^*, \sigma^*) \leftarrow \mathcal{A}^{\Sign(\textit{sk}, \cdot), \mathcal{H}(\cdot)}(\textit{pk})
    \end{array}
  \right] \leq \epsilon,
\]
where $\mu^*$ was not queried to the signing oracle.
\end{definition}

\subsection{Main Theorem}

\begin{theorem}[EUF-CMA Security of TL-DSA]
\label{thm:main}
In the random oracle model, for any EUF-CMA adversary $\mathcal{A}$ against TL-DSA making at most $q_S$ signing queries and $q_H$ random oracle queries, there exists an MSIS adversary $\mathcal{B}$ such that:
\[
  \Adv^{\EUFCMA}_{\textsf{TL-DSA}}(\mathcal{A}) \leq \frac{q_H + q_S + 1}{3^{243}} + \Adv^{\MSIS}_{k, l+1, 1}(\mathcal{B}),
\]
where the MSIS instance has parameters $(k, l+1, \beta=1)$ over $R_3$ with $n = 256$.
The running time of $\mathcal{B}$ is that of $\mathcal{A}$ plus $O(q_H \cdot k \cdot l \cdot n^2)$ for polynomial arithmetic.
\end{theorem}

\subsection{Proof of Theorem~\ref{thm:main}}

The proof follows the standard Fiat--Shamir signature security framework, adapted to exploit the ternary non-abort property.

\medskip
\noindent\textbf{Game 0.} This is the real EUF-CMA experiment. $\mathcal{A}$ receives a public key $(\rho, \mathbf{t})$ where $\mathbf{t} = \mathbf{A} \cdot \mathbf{s}_1$, has access to signing oracle $\Sign(\textit{sk}, \cdot)$ and random oracle $\mathcal{H}(\cdot)$, and outputs a forgery attempt $(\mu^*, \sigma^* = (\mathbf{z}^*, \tilde{c}^*))$.

\medskip
\noindent\textbf{Game 1.} We replace the random oracle $\mathcal{H}$ with a lazily-sampled table, answering each fresh query with a uniformly random element of $\{-1,0,+1\}^{243}$. This is a purely syntactic change:
\[
  \Pr[\text{Game 1 succeeds}] = \Pr[\text{Game 0 succeeds}].
\]

\medskip
\noindent\textbf{Game 2.} We simulate the signing oracle without using $\mathbf{s}_1$. By Theorem~\ref{thm:zk}, signatures can be perfectly simulated:

For each signing query on message $\mu_i$:
\begin{enumerate}[leftmargin=2em]
\item Sample $\mathbf{z}_i \xleftarrow{\$} R_3^l$.
\item Compute $\bar{\mu}_i = \mathcal{H}(\texttt{trits}(\mathbf{t}) \| \mu_i)$.
\item Sample $\tilde{c}_i \xleftarrow{\$} \{-1,0,+1\}^{243}$.
\item Compute $c_i = \textsf{SampleChallenge}(\tilde{c}_i)$.
\item Compute $\mathbf{w}_i = \mathbf{A} \cdot \mathbf{z}_i - c_i \cdot \mathbf{t}$.
\item Program $\mathcal{H}(\bar{\mu}_i \| \texttt{trits}(\mathbf{w}_i)) := \tilde{c}_i$.
\end{enumerate}

Programming fails only if $(\bar{\mu}_i \| \texttt{trits}(\mathbf{w}_i))$ was previously queried to $\mathcal{H}$, which happens with probability at most $(q_H + q_S) / 3^{243}$ per query (since $\mathbf{w}_i$ depends on the freshly sampled $\mathbf{z}_i$). By a union bound over $q_S$ signing queries:
\[
  \left| \Pr[\text{Game 2}] - \Pr[\text{Game 1}] \right| \leq \frac{q_S(q_H + q_S)}{3^{243}}.
\]

Since $3^{243}$ is astronomically large ($\approx 2^{385}$), this term is negligible for any polynomial $q_S, q_H$.

\medskip
\noindent\textbf{Game 2 analysis: extracting an MSIS solution.}

In Game~2, the signing oracle does not use $\mathbf{s}_1$. Therefore, $\mathcal{A}$'s view depends on $\mathbf{s}_1$ only through the public key $\mathbf{t} = \mathbf{A} \cdot \mathbf{s}_1$. Suppose $\mathcal{A}$ outputs a valid forgery $(\mu^*, \mathbf{z}^*, \tilde{c}^*)$ where $\Verify(\textit{pk}, \mu^*, (\mathbf{z}^*, \tilde{c}^*)) = 1$. Then:
\begin{align}
  \mathbf{A} \cdot \mathbf{z}^* - c^* \cdot \mathbf{t} &= \mathbf{w}^*, \label{eq:verify1}\\
  \mathcal{H}(\bar{\mu}^* \| \texttt{trits}(\mathbf{w}^*)) &= \tilde{c}^*, \label{eq:verify2}
\end{align}
where $c^* = \textsf{SampleChallenge}(\tilde{c}^*)$.

Since $\mathbf{t} = \mathbf{A} \cdot \mathbf{s}_1$, equation~\eqref{eq:verify1} gives:
\[
  \mathbf{A} \cdot \mathbf{z}^* - c^* \cdot \mathbf{A} \cdot \mathbf{s}_1 = \mathbf{w}^* \quad \Longrightarrow \quad \mathbf{A} \cdot (\mathbf{z}^* - c^* \cdot \mathbf{s}_1) = \mathbf{w}^*.
\]

Now we apply the \textbf{forking lemma}. We rewind $\mathcal{A}$ to the point where $\tilde{c}^*$ was programmed and provide a different random oracle response $\tilde{c}^{**} \neq \tilde{c}^*$. With probability at least $1/(q_H + q_S + 1)$ (by the general forking lemma), $\mathcal{A}$ produces a second valid forgery $(\mu^*, \mathbf{z}^{**}, \tilde{c}^{**})$ on the same message with the same $\mathbf{w}^*$ but different challenge $c^{**} = \textsf{SampleChallenge}(\tilde{c}^{**})$.

From the two forgeries:
\begin{align*}
  \mathbf{A} \cdot \mathbf{z}^* - c^* \cdot \mathbf{t} &= \mathbf{w}^*,\\
  \mathbf{A} \cdot \mathbf{z}^{**} - c^{**} \cdot \mathbf{t} &= \mathbf{w}^*.
\end{align*}
Subtracting:
\[
  \mathbf{A} \cdot (\mathbf{z}^* - \mathbf{z}^{**}) = (c^* - c^{**}) \cdot \mathbf{t} = (c^* - c^{**}) \cdot \mathbf{A} \cdot \mathbf{s}_1.
\]
Therefore:
\[
  \mathbf{A} \cdot \left[ (\mathbf{z}^* - \mathbf{z}^{**}) - (c^* - c^{**}) \cdot \mathbf{s}_1 \right] = \mathbf{0}.
\]

Define the vector $\mathbf{v} = \begin{pmatrix} \mathbf{z}^* - \mathbf{z}^{**} \\ -(c^* - c^{**}) \end{pmatrix} \in R_3^{l+1}$, and extend the matrix as $\mathbf{A}' = [\mathbf{A} \mid \mathbf{t}] \in R_3^{k \times (l+1)}$. Then:
\[
  \mathbf{A}' \cdot \mathbf{v} = \mathbf{0}, \quad \linf{\mathbf{v}} \leq 1, \quad \mathbf{v} \neq \mathbf{0} \text{ (since } c^* \neq c^{**}\text{)}.
\]

This is a valid solution to the $\MSIS_{k, l+1, 1}$ problem on the matrix $\mathbf{A}'$.

\medskip
\noindent\textbf{Combining the bounds.} We construct $\mathcal{B}$ as follows: $\mathcal{B}$ receives an MSIS challenge $\mathbf{A}'$, parses it as $[\mathbf{A} \mid \mathbf{t}]$, runs $\mathcal{A}$ in Game~2 using $(\rho, \mathbf{t})$ as the public key (where $\rho$ is set to reproduce $\mathbf{A}$), and applies the forking lemma to extract $\mathbf{v}$.

\[
  \Adv^{\EUFCMA}_{\textsf{TL-DSA}}(\mathcal{A})
  \leq \frac{q_S(q_H + q_S)}{3^{243}} + (q_H + q_S + 1) \cdot \Adv^{\MSIS}_{k, l+1, 1}(\mathcal{B}).
\]

Since $3^{243} \approx 2^{385}$, the first term is negligible for any polynomial number of queries. \qed

% ===================================================================
% SECTION 6: MODULE-SIS TO MODULE-LWE DUALITY
% ===================================================================
\section{Module-SIS to Module-LWE Duality}
\label{sec:duality}

The hardness of MSIS over $R_3$ can be related to the hardness of MLWE over $R_3$ via the standard lattice duality results.

\begin{theorem}[SIS--LWE duality, adapted from {\cite[Theorem 4.1]{peikert2016decade}}]
\label{thm:duality}
For the ring $R_3 = \Z_3[X]/(X^{256}+1)$, solving $\MSIS_{k, l, \beta}$ with $\beta = 1$ is at least as hard as solving $\MLWE_{k, \chi}$ where $\chi$ is any distribution over $R_3$. Formally, for any MSIS solver $\mathcal{B}$, there exists an MLWE distinguisher $\mathcal{D}$ such that:
\[
  \Adv^{\MSIS}_{k, l, 1}(\mathcal{B}) \leq \Adv^{\MLWE}_{k, \chi}(\mathcal{D}) + \frac{1}{3^{nk}}.
\]
\end{theorem}

\begin{proof}[Proof sketch]
The proof follows the standard approach. Given an MLWE instance $(\mathbf{A}, \mathbf{b})$ where either $\mathbf{b} = \mathbf{A} \cdot \mathbf{s} + \mathbf{e}$ or $\mathbf{b} \xleftarrow{\$} R_3^k$:

If $\mathbf{b}$ is uniform, then $[\mathbf{A} \mid \mathbf{b}]$ is a uniform matrix in $R_3^{k \times (k+1)}$, and any MSIS solution on this matrix would solve a uniformly random MSIS instance.

If $\mathbf{b} = \mathbf{A} \cdot \mathbf{s} + \mathbf{e}$, then $[\mathbf{A} \mid \mathbf{b}]$ has the hidden short vector $(\mathbf{s}^\top \mid -1)^\top$ satisfying $[\mathbf{A} \mid \mathbf{b}] \cdot (\mathbf{s}^\top \mid -1)^\top = \mathbf{e}$ (which is short). An MSIS solver that finds \emph{any} short vector in the kernel would either find this planted solution or an independent one---either way, the existence of a short kernel vector distinguishes the two cases.

Over $R_3$, the key observation is that $\beta = 1$ and all elements already satisfy $\linf{\cdot} \leq 1$, so the ``shortness'' condition is automatically satisfied. The hardness therefore reduces to finding \emph{any} nontrivial kernel element of a random matrix over $R_3$, which is the core MLWE/MSIS problem in the ternary setting.
\end{proof}

\begin{corollary}
The EUF-CMA security of TL-DSA reduces to the hardness of Module-LWE over $R_3$:
\[
  \Adv^{\EUFCMA}_{\textsf{TL-DSA}}(\mathcal{A})
  \leq \frac{q_S(q_H + q_S)}{3^{243}} + (q_H + q_S + 1) \cdot \left( \Adv^{\MLWE}_{k, \chi}(\mathcal{D}) + \frac{1}{3^{nk}} \right).
\]
\end{corollary}

% ===================================================================
% SECTION 7: CONCRETE SECURITY ANALYSIS
% ===================================================================
\section{Concrete Security Estimates}
\label{sec:concrete}

\subsection{Lattice Dimension and Hardness}

The concrete security of TL-DSA depends on the hardness of finding short vectors in Module-LWE/SIS lattices over $R_3$.

For a Module-LWE instance with matrix $\mathbf{A} \in R_3^{k \times l}$ and polynomial degree $n = 256$, the underlying lattice has rank $d = (k + l) \cdot n$ over $\Z_3$. The Hermite factor $\delta$ required to solve this lattice problem using BKZ with block size $b$ satisfies:
\[
  \delta^d \approx \frac{3^{nk/d}}{\sqrt{2\pi e}},
\]
and the best known classical algorithms achieve $\delta \approx \left( \frac{\pi b}{2\pi e} \cdot (\pi b)^{1/b} \right)^{1/(2(b-1))}$.

\begin{table}[ht]
\centering
\caption{Concrete security estimates for TL-DSA parameter sets.}
\label{tab:security}
\begin{tabular}{@{}lcccccc@{}}
\toprule
Variant & $k$ & $l$ & Lattice rank $d$ & Estimated BKZ-$b$ & Classical bits & Quantum bits \\
\midrule
TL-DSA-44 & 4 & 4 & 2048 & $\geq 356$ & $\geq 128$ & $\geq 116$ \\
TL-DSA-65 & 6 & 5 & 2816 & $\geq 538$ & $\geq 192$ & $\geq 174$ \\
TL-DSA-87 & 8 & 7 & 3840 & $\geq 720$ & $\geq 256$ & $\geq 232$ \\
\bottomrule
\end{tabular}
\end{table}

The estimates use the Core-SVP methodology: classical security is $0.292b$ bits (sieving) and quantum security is $0.265b$ bits (Grover-accelerated sieving), consistent with the methodology used in the ML-DSA specification.

\subsection{Hash Function Security}

The ternary sponge has capacity 486 trits ($= 3^{486} \approx 2^{770}$ states). Collision resistance is bounded by the birthday attack at $\sqrt{3^{486}} \approx 2^{385}$ operations. Preimage resistance is $3^{486} \approx 2^{770}$ operations. Both are superexponential in the security parameter for all TL-DSA variants.

The challenge space has $\binom{256}{\tau} \cdot 2^\tau$ elements (choosing $\tau$ positions from 256 and assigning each $\pm 1$):

\begin{table}[ht]
\centering
\caption{Challenge space sizes.}
\label{tab:challenge}
\begin{tabular}{@{}lccc@{}}
\toprule
Variant & $\tau$ & $\log_2\binom{256}{\tau} \cdot 2^\tau$ & Collision security \\
\midrule
TL-DSA-44 & 39 & $\approx 178$ & $\geq 89$ bits \\
TL-DSA-65 & 49 & $\approx 213$ & $\geq 106$ bits \\
TL-DSA-87 & 60 & $\approx 248$ & $\geq 124$ bits \\
\bottomrule
\end{tabular}
\end{table}

The challenge collision security is the binding constraint for TL-DSA-44 and TL-DSA-65. For TL-DSA-87, the lattice dimension is the binding constraint.

\subsection{Sizes}

\begin{table}[ht]
\centering
\caption{Key and signature sizes (in trits and bytes).}
\label{tab:sizes}
\begin{tabular}{@{}lcccccc@{}}
\toprule
Variant & $|\textit{pk}|$ (trits) & $|\textit{pk}|$ (bytes) & $|\textit{sk}|$ (trits) & $|\textit{sk}|$ (bytes) & $|\sigma|$ (trits) & $|\sigma|$ (bytes) \\
\midrule
TL-DSA-44 & 1267 & 317 & 3315 & 829 & 1267 & 317 \\
TL-DSA-65 & 1779 & 445 & 4595 & 1149 & 1523 & 381 \\
TL-DSA-87 & 2291 & 573 & 6131 & 1533 & 2035 & 509 \\
\bottomrule
\end{tabular}
\end{table}

Byte sizes computed as $\lceil \text{trits} \times 2 / 8 \rceil$ (2 bits per trit in packed encoding, matching the \texttt{to\_bytes} implementation in \texttt{ternary\_lattice.rs} line~217).

% ===================================================================
% SECTION 8: DISCUSSION
% ===================================================================
\section{Discussion}
\label{sec:discussion}

\subsection{Comparison to ML-DSA}

TL-DSA and ML-DSA share the same high-level structure (Fiat--Shamir with Aborts over Module-LWE lattices) but differ in three significant ways:

\begin{enumerate}[leftmargin=2em]
\item \textbf{Ring modulus.} ML-DSA uses $q = 8380417$ with NTT-friendly factorization. TL-DSA uses $q = 3$, which precludes NTT but enables the non-abort property. Ring multiplication is $O(n^2)$ schoolbook rather than $O(n \log n)$ NTT.

\item \textbf{Rejection sampling.} ML-DSA requires $\sim$4--7 signing attempts on average. TL-DSA requires exactly one (Theorem~\ref{thm:nonabort}), yielding deterministic signing with perfect zero-knowledge.

\item \textbf{Key and signature sizes.} TL-DSA produces significantly smaller keys and signatures because each coefficient is a trit (2 bits packed) rather than a $\lceil \log_2 q \rceil = 23$-bit integer. However, the security per trit of lattice dimension is lower because the lattice is over a smaller ring.
\end{enumerate}

\subsection{Security Considerations Specific to $q=3$}

The small modulus $q=3$ has implications that require careful analysis:

\textbf{Algebraic attacks.} Over $\Z_3[X]/(X^{256}+1)$, the polynomial $X^{256}+1$ factors nontrivially modulo 3 (since $X^{256}+1 = (X^{128})^2 + 1$ and $-1 \equiv 2 \pmod{3}$). The factorization structure could potentially be exploited by algebraic attacks that do not apply to NTT-friendly primes. The security estimates in Table~\ref{tab:security} assume generic lattice attacks; algebraic speedups are an open research question.

\textbf{Entropic hardness.} When $q=3$, the secret $\mathbf{s}_1$ and the uniform distribution have the same support ($\{-1,0,+1\}$). The MLWE problem reduces to distinguishing $\mathbf{A} \cdot \mathbf{s}$ from uniform, which requires exploiting the algebraic structure of $\mathbf{A}$ rather than the ``smallness'' of $\mathbf{s}$. The module rank $(k, l)$ must therefore be sufficiently large to compensate for the lack of a noise-to-modulus ratio.

\subsection{Open Problems}

\begin{enumerate}[leftmargin=2em]
\item \textbf{Tight reduction.} The forking lemma introduces a $q_H$-factor loss in the reduction. Achieving a tight reduction (without the forking lemma) via the lossy identification paradigm is an open question for TL-DSA.

\item \textbf{Algebraic factorization analysis.} A complete analysis of the factorization of $X^{256}+1$ over $\Z_3$ and its impact on lattice attacks would strengthen the concrete security estimates.

\item \textbf{Formal security of the ternary sponge.} The proof assumes the ternary sponge is a random oracle. A dedicated cryptanalysis establishing the sponge's indifferentiability from a random oracle would close this gap (see companion report: Ternary Sponge Cryptanalysis, Task~2.5).
\end{enumerate}

% ===================================================================
% SECTION 9: CONCLUSION
% ===================================================================
\section{Conclusion}

We have established the EUF-CMA security of TL-DSA under the Module-LWE hardness assumption over the ternary polynomial ring $R_3 = \Z_3[X]/(X^{256}+1)$, in the random oracle model. The security reduction proceeds via a standard Fiat--Shamir argument with forking, with the key innovation that the ternary non-abort property (Theorem~\ref{thm:nonabort}) eliminates rejection sampling entirely, yielding deterministic signing with perfect zero-knowledge simulation. Concrete security estimates confirm that the three parameter sets (TL-DSA-44/65/87) achieve their targeted 128/192/256-bit classical security levels under standard lattice attack cost models.

% ===================================================================
% REFERENCES
% ===================================================================
\begin{thebibliography}{99}

\bibitem{dilithium}
L.~Ducas, E.~Kiltz, T.~Lepoint, V.~Lyubashevsky, P.~Schwabe, G.~Seiler, and D.~Stehl\'{e}.
\newblock CRYSTALS--Dilithium: A Lattice-Based Digital Signature Scheme.
\newblock \textit{IACR Transactions on Cryptographic Hardware and Embedded Systems}, 2018(1):238--268, 2018.

\bibitem{lyubashevsky2012}
V.~Lyubashevsky.
\newblock Lattice Signatures without Trapdoors.
\newblock In \textit{EUROCRYPT 2012}, LNCS 7237, pages 738--755. Springer, 2012.

\bibitem{fiat-shamir}
A.~Fiat and A.~Shamir.
\newblock How to Prove Yourself: Practical Solutions to Identification and Signature Problems.
\newblock In \textit{CRYPTO 1986}, LNCS 263, pages 186--194. Springer, 1986.

\bibitem{fips204}
National Institute of Standards and Technology.
\newblock FIPS 204: Module-Lattice-Based Digital Signature Standard (ML-DSA).
\newblock August 2024.

\bibitem{peikert2016decade}
C.~Peikert.
\newblock A Decade of Lattice Cryptography.
\newblock \textit{Foundations and Trends in Theoretical Computer Science}, 10(4):283--424, 2016.

\bibitem{langlois-stehle}
A.~Langlois and D.~Stehl\'{e}.
\newblock Worst-Case to Average-Case Reductions for Module Lattices.
\newblock \textit{Designs, Codes and Cryptography}, 75(3):565--599, 2015.

\bibitem{kiltz-loss}
E.~Kiltz, V.~Lyubashevsky, and C.~Schaffner.
\newblock A Concrete Treatment of Fiat-Shamir Signatures in the Quantum Random Oracle Model.
\newblock In \textit{EUROCRYPT 2018}, LNCS 10822, pages 552--586. Springer, 2018.

\bibitem{forking}
D.~Pointcheval and J.~Stern.
\newblock Security Arguments for Digital Signatures and Blind Signatures.
\newblock \textit{Journal of Cryptology}, 13(3):361--396, 2000.

\end{thebibliography}

% ===================================================================
% APPENDIX
% ===================================================================
\appendix

\section{Implementation Correspondence}
\label{app:impl}

The following table maps each algorithmic step to the production Rust implementation in the PlenumNET repository.

\begin{table}[ht]
\centering
\small
\caption{Algorithm-to-code correspondence.}
\begin{tabular}{@{}lll@{}}
\toprule
Algorithm Step & Source File & Lines \\
\midrule
$R_3$ ring definition & \texttt{ternary\_lattice.rs} & 58--62, 95--99 \\
\texttt{mod3} reduction & \texttt{ternary\_lattice.rs} & 67--75 \\
Ring multiplication (schoolbook) & \texttt{ternary\_lattice.rs} & 185--214 \\
CBD$_\eta$ sampling & \texttt{ternary\_lattice.rs} & 423--449 \\
Uniform ternary sampling & \texttt{ternary\_lattice.rs} & 401--421 \\
Matrix expansion & \texttt{ternary\_lattice.rs} & 452--461 \\
Parameter sets & \texttt{tl\_dsa.rs} & 73--103 \\
$\KeyGen$ & \texttt{tl\_dsa.rs} & 198--231 \\
$\Sign$ (with abort loop) & \texttt{tl\_dsa.rs} & 254--308 \\
$\Verify$ & \texttt{tl\_dsa.rs} & 310--348 \\
$\textsf{SampleChallenge}$ & \texttt{tl\_dsa.rs} & 157--183 \\
Ternary sponge hash & \texttt{sponge.rs} & 46--69 (permutation) \\
$\ell_\infty$ norm & \texttt{ternary\_lattice.rs} & 141--143 \\
\bottomrule
\end{tabular}
\end{table}

\section{Test Vector Correspondence}
\label{app:tests}

The following unit tests in \texttt{tl\_dsa.rs} validate the implementation against the formal specification:

\begin{itemize}[leftmargin=2em]
\item \texttt{test\_keygen\_deterministic} (line~422): Verifies that $\KeyGen$ is deterministic given the same seed.
\item \texttt{test\_sign\_verify\_44} (line~438): End-to-end correctness for TL-DSA-44.
\item \texttt{test\_sign\_verify\_wrong\_message} (line~449): Forgery detection (wrong message).
\item \texttt{test\_sign\_verify\_wrong\_key} (line~462): Forgery detection (wrong key).
\item \texttt{test\_sign\_deterministic} (line~478): Verifies deterministic signing (Corollary~\ref{cor:deterministic}).
\item \texttt{test\_sample\_challenge} (line~514): Verifies challenge weight bound $\wt(c) \leq \tau$.
\end{itemize}

All tests pass as of the March 2026 audit.

\end{document}